\documentclass[11pt]{article}
\usepackage{graphicx} % Required for inserting images
\usepackage{siunitx}
\usepackage{esint}
\usepackage{enumitem}
\usepackage{amsmath}
\usepackage{tikz-feynman}
\usepackage{float}

\newcommand{\lagr}{\mathcal{L}}
\newcommand{\half}{\frac{1}{2}}
\DeclareSIUnit{\b}{b}
\DeclareSIUnit{\pb}{pb}

\usepackage{geometry}
 \geometry{
 total={170mm,257mm},
 left=16mm,
 top=24mm,
 bottom=30mm,
 top=15mm
 }

\title{Intro to Particles \\ HW 3}
\author{Erel Zack \\ erel.zack@campus.technion.ac.il \and Benjamin Leshansky
 \\ benjaminl@campus.technion.ac.il}
\date{\today}

\begin{document}
\maketitle

\section{Symmetry to Rotations}

\subsection{}

\begin{equation}
\begin{split}
    [L_i,L_j] &= 
    [\varepsilon_{imn}x_m p_n,\, \varepsilon_{jkl}x_k p_l] =
    \varepsilon_{imn}\varepsilon_{jkl}[x_mp_n, x_kp_l]=\\
    &=
    \varepsilon_{imn}\varepsilon_{jkl} (x_m[p_n,x_kp_l] + [x_m,x_kp_l]p_n)=
    \varepsilon_{imn}\varepsilon_{jkl}(x_m[p_n,x_k]p_l + x_k[x_m,p_l]p_n)=\\
    &= 
    i\varepsilon_{imn}\varepsilon_{jkl}(-\delta_{nk}x_m p_l + \delta_{lm}x_k p_n) = 
    -i\varepsilon_{imk}\varepsilon_{jkl}x_mp_l + i\varepsilon_{imn}\varepsilon_{jkm}x_kp_n=\\
    &=
    i\varepsilon_{imk}\varepsilon_{jlk}x_m p_l - i\varepsilon_{inm}\varepsilon_{jkm}x_k p_n =
    i(\delta_{ij}\delta_{ml}-\delta_{il}\delta_{mj})x_m p_l - i(\delta_{ij}\delta_{nk}-\delta_{ik}\delta_{nj})x_k p_n =\\
    &=
    i\delta_{ij}x_m p_m - i x_j p_i - i\delta_{ij} x_n p_n + i x_i p_j =
    i\delta_{ij}(\vec x\cdot\vec p) + i(x_ip_j-x_jp_i) - i\delta_{ij}(\vec x\cdot\vec p) =\\
    &=
    i(x_i p_j-x_j p_i)=
    i\varepsilon_{ijk}L_k
\end{split}
\end{equation}

\subsection{}
The rotational transformation is $g(\theta , \phi , \psi ) = R_1 (\theta) R_2(\phi)R_3(\psi)$, thus we can use write explicitly the inf. repr. here: $$$$ 

\begin{equation}
\begin{split}
        &g(\theta, \phi, \psi ) = \left(I+\begin{pmatrix} 0&0&0 \\ 0 & 0 & -\theta \\ 0&\theta&0 \end{pmatrix} \right)\left(I+\begin{pmatrix} 0&0&-\phi \\ 0 & 0 &0 \\ \phi&0&0 \end{pmatrix} \right)\left(I+\begin{pmatrix} 0&-\psi&0 \\ \psi & 0 & 0 \\ 0&0&0 \end{pmatrix} \right) + \mathcal{O}(\theta^2,\phi^2,\psi^2) \\ & = 
        I + \begin{pmatrix} 0&0&0 \\ 0 & 0 & -\theta \\ 0&\theta&0 \end{pmatrix} + \begin{pmatrix} 0&0&-\phi \\ 0 & 0 &0 \\ \phi&0&0 \end{pmatrix} + \begin{pmatrix} 0&-\psi&0 \\ \psi & 0 & 0 \\ 0&0&0 \end{pmatrix} + \mathcal{O}(\theta^2,\phi^2,\psi^2)
\end{split}
\end{equation}

This is the inf. representation of the rotation group


\subsection{}

Approximating for infinitesimal rotations:
\begin{equation}
    e^{-i\vec \alpha\cdot \vec L} \simeq 
    I - i\vec \alpha\cdot\vec L =
    I-i\theta L_x -i\varphi L_y - i\psi L_z
\end{equation}
Comparing to the transformation, we equate the coefficients for the euler angles:
\begin{equation}
    L_x = i
    \begin{pmatrix}
        0&0&0\\
        0&0&-\theta\\
        0&\theta&0
    \end{pmatrix}
    ,\quad
    L_y = i
    \begin{pmatrix}
        0&0&\phi\\
        0&0&0\\
        \phi&0&0
    \end{pmatrix}
    ,\quad
    L_z = i
    \begin{pmatrix}
        0&-\psi&0\\
        \psi&0&0\\
        0&0&0
    \end{pmatrix}
\end{equation}


\section{Constructing Dirac Lagrangian}
\subsection{}
Using the hint:
\begin{equation}
\begin{split}
    [\mathcal S] &= \left[\int\mathcal L\,d^4x \right]\\
    1 &= [\mathcal L]\cdot [x]^4\\
    1 &= [\mathcal L]\cdot M^{-4}\\
    \Rightarrow [\mathcal L] &= M^4
\end{split}
\end{equation}

\subsection{}
Using our previous result for $\mathcal L$:
\begin{equation}
\begin{split}
    [\mathcal L] &= \left[\half \partial_\mu\phi\partial^\mu\phi \right]\\
    M^4 &= [\partial_\mu]^2[\phi]^2\\
    M^4 &= \frac{[\phi]^2}{[x]^2}\\
    M^4 &= [\phi]^2 M^2\\
    [\phi]^2 &= M^2\\
    [\phi] &= M
\end{split}
\end{equation}

\subsection{}

\begin{equation}
    \begin{split}
        S^{0i}=\frac{i}{4}[\gamma^0,\gamma^i] = 
        \frac{i}{4}(-\{ \gamma^0,\gamma^i\} +2\gamma^0\gamma^i)=
        \frac{i}{4}(-2g^{0i} + 2\gamma^0\gamma^i)=
        \frac{i}{2}\gamma^0\gamma^i
    \end{split}
\end{equation}

Note that by its definition:
\begin{equation}
\begin{split}
    \gamma^i&=\beta a^i\\
    \Rightarrow \beta \gamma^i &= \beta^2 a^i = a^i\\
    \Rightarrow \gamma^0\gamma^i &= a^i
\end{split}
\end{equation}

Which gives us:
\begin{equation}
    S^{0i} = \frac{i}{2}a^i =
    \frac{i}{2}
    \begin{pmatrix}
        0&\sigma^i\\
        \sigma^i & 0
    \end{pmatrix}
\end{equation}

Now: $$S^{ij} = \frac{i}{4}[\gamma^i,\gamma^j] = \frac{i}{4}(\{\gamma^i,\gamma^j\} - 2\gamma^j\gamma^i) = \frac{i}{2}(g^{ij}-\gamma^j\gamma^i) = -\frac{i}{2}(I+\gamma^j\gamma^i)$$


\subsection{}
We now develop in TAYLOR and take the first order:
\begin{equation}
    \psi\rightarrow\psi ' = \Lambda_{\half} \psi = \left(I - \frac{i}{2}\omega_{\mu\nu}S^{\mu\nu}\right)\psi
\end{equation}

\noindent
And for $\bar\psi$:
$$\bar\psi=\psi^\dagger \gamma^0 \to (\Lambda \psi)^\dagger \gamma^0 = \psi^\dagger \Lambda^\dagger \gamma^0=(\psi^\dagger \gamma^0)\gamma^0\Lambda^\dagger \gamma^0 = \bar\psi\gamma^0\Lambda^\dagger\gamma^0$$

$$\gamma^0S^{\mu \nu}\gamma^0 = \frac{i}{4}(\gamma^0 \gamma^\mu \gamma^\nu \gamma^0 - \gamma^0 \gamma^\nu \gamma^\mu \gamma^0)$$

Calculating separately for the $0i$ and $ij$ cases:
$$\gamma^0S^{0i}\gamma^0 = \frac{i}{4}(\gamma^0 \gamma^0 \gamma^i \gamma^0 - \gamma^0 \gamma^i \gamma^0 \gamma^0) = \frac{i}{4}(\gamma^i\gamma^0-\gamma^0\gamma^i)=\frac{i}{4}[\gamma^i,\gamma^0]= -S^{0i}$$

As we've shown before:
\begin{equation}
\begin{split}
    &S^{0i}=\frac{i}{4}[\gamma^0,\gamma^i]= \frac{i}{2}\gamma^0\gamma^i\\
    \Rightarrow &[\gamma^0,\gamma^i]=2\gamma^0\gamma^i\\
    &\gamma^0\gamma^i-\gamma^i\gamma^0 = 2\gamma^0\gamma^i\\
    &\gamma^0\gamma^i = -\gamma^i\gamma^0
\end{split}
\end{equation}

And therefore:
\begin{equation}
\begin{split}
    \gamma^0 S^{ij}\gamma^0 = 
    -\frac{i}{2}\gamma^0(I+\gamma^j\gamma^i)\gamma^0=
    -\frac{i}{2}(I+\gamma^0\gamma^j\gamma^i\gamma^0)=
    -\frac{i}{2}(I+(-\gamma^j\gamma^0)(-\gamma^0\gamma^i))=
    S^{ij}
\end{split}
\end{equation}

This gives us:
\begin{equation}
\begin{split}
    &\gamma^0\Lambda^\dagger\gamma^0 = \gamma^0\left(I-\frac{i}{2}\omega_{\mu\nu}S^{\mu\nu} \right)\gamma^0 = 
    I-\frac{i}{2}\omega_{\mu\nu}\gamma^0 S^{\mu\nu}\gamma^0 =\\
    &=I-\frac{i}{2}\left(-\omega_{0i}S^{0i} + \omega_{ij}S^{ij}  \right) = I - \frac{i}{2}( - \omega_{\mu\nu}S^{\mu\nu})=I+ \frac{i}{2}\omega_{\mu\nu}S^{\mu\nu}
\end{split}
\end{equation}
Now: 
\begin{equation}
    \bar\psi\psi\to \bar\psi\left(I+\frac{i}{2}\omega_{\mu\nu}S^{\mu\nu} \right)\left(I-\frac{i}{2}\omega_{\mu\nu}S^{\mu\nu} \right)\psi = 
    \bar\psi\left( I + \left(\frac{1}{2}\omega_{\mu\nu}S^{\mu\nu} \right)^2\right)\psi = \bar\psi\psi + \mathcal O(\omega_{\mu\nu}^2)
\end{equation}

Meanwhile:
$$\psi^\dagger\psi\to (\Lambda\psi)^\dagger(\Lambda\psi)=\psi^\dagger\Lambda^\dagger\Lambda\psi  $$
But 
$$\Lambda^\dagger=(1-\frac{i}{2}\omega_{\mu\nu} S^{\mu\nu})^\dagger=1+\frac{i}{2}\omega_{\mu\nu}(S^{\mu\nu})^\dagger =1 + \frac{i}{2}(\omega_{ij}S^{ij}-\omega_{0i}S^{0i})$$

Where we used:
\begin{equation}
\begin{split}
    &(S^{\mu\nu})^\dagger = 
    -\frac{i}{4}[\gamma^\mu,\gamma^\nu]^\dagger=
    -\frac{i}{4}[\gamma^{\nu\dagger},\gamma^{\mu\dagger}]=
    \frac{i}{4}[\gamma^{\mu\dagger},\gamma^{\nu\dagger}]\\
    \Rightarrow & (S^{ij})^\dagger = 
    \frac{i}{4}[\gamma^{i\dagger},\gamma^{j\dagger}]=
    \frac{i}{4}[-\gamma^i,-\gamma^j]=
    \frac{i}{4}[\gamma^i,\gamma^j] = 
    S^{ij}\\
    \Rightarrow & (S^{0i})^\dagger = 
    \frac{i}{4}[\gamma^{0\dagger},\gamma^{i\dagger}]=
    \frac{i}{4}[\gamma^0,-\gamma^i]=
    -S^{0i}
\end{split}
\end{equation}

This gives us:
\begin{equation}
\begin{split}
    \psi^\dagger\Lambda^\dagger\Lambda\psi &= 
    \psi^\dagger \left( I + \frac{i}{2}(\omega_{ij}S^{ij}-\omega_{0i}S^{0i}) \right) \left(I - \frac{i}{2}(\omega_{ij}S^{ij} + \omega_{0i}S^{0i}) \right) \psi=\\
    &=\psi^\dagger\left(I - i\omega_{0i}S^{0i} +\mathcal O(\omega_{\mu \nu}^2) \right)\psi  \neq \psi^\dagger\psi + \mathcal O(\omega^2)
\end{split}
\end{equation}

\subsection{}
We know that $[-m\bar\psi \psi]=[\mathcal L]=M^4$ therefore: $$ [\bar\psi \psi ] = M^3$$ thus \begin{equation}
    \boxed{[\psi]=M^{3/2}}
\end{equation}

Therefore the kinetic term will be of the form:
\begin{equation}
\begin{split}
    [\partial_\mu]^n[\psi]^2 &= [\mathcal L]\\
    M^n\cdot M^3 &= M^4\\
    M^n &= M\\
    n&=1
\end{split}
\end{equation}

\subsection{}

Define $$\Lambda^\mu{}_\nu = e^{-\frac{i}{2}\omega_{\rho\sigma}(\mathcal J^{\rho\sigma})^\mu{}_\nu}$$

\subsubsection{}

\[
S^{\rho\sigma}=\frac{i}{4}[\gamma^\rho,\gamma^\sigma]
\]

\[
(\mathcal J^{\rho\sigma})^\mu{}_\nu\gamma^\nu
=
[\gamma^\mu,S^{\rho\sigma}]
\]

\[
(\mathcal J^{\rho\sigma})^\mu{}_\nu
=
i g^{\mu\alpha}
\left(
\delta^\rho_\alpha\delta^\sigma_\nu
-
\delta^\sigma_\alpha\delta^\rho_\nu
\right)
\]

\[
(\mathcal J^{\rho\sigma})^\mu{}_\nu
=
i
\left(
g^{\mu\rho}\delta^\sigma_\nu
-
g^{\mu\sigma}\delta^\rho_\nu
\right)
\]

\[
(\mathcal J^{\rho\sigma})^\mu{}_\nu\gamma^\nu
=
i
\left(
g^{\mu\rho}\gamma^\sigma
-
g^{\mu\sigma}\gamma^\rho
\right)
\]

\[
[\gamma^\mu,S^{\rho\sigma}]
=
\frac{i}{4}
[\gamma^\mu,[\gamma^\rho,\gamma^\sigma]]
\]

\[
[\gamma^\rho,\gamma^\sigma]
=
\gamma^\rho\gamma^\sigma-\gamma^\sigma\gamma^\rho
\]
\[
[\gamma^\mu,S^{\rho\sigma}]
=
\frac{i}{4}
\left(
[\gamma^\mu,\gamma^\rho\gamma^\sigma]
-
[\gamma^\mu,\gamma^\sigma\gamma^\rho]
\right)
\]
\[
\{\gamma^\mu,\gamma^\nu\}=2g^{\mu\nu}
\]
\[
[\gamma^\mu,\gamma^\rho\gamma^\sigma]
= (2g^{\mu\rho}-\gamma^\rho\gamma^\mu)\gamma^\sigma - \gamma^\rho \gamma^\sigma \gamma^\mu = 2g^{\mu\rho}\gamma^\sigma
-
\gamma^\rho(\gamma^\sigma\gamma^\mu+\gamma^\mu\gamma^\sigma) = 
2g^{\mu\rho}\gamma^\sigma
-
\gamma^\rho2g^{\mu\sigma}
\]
\[
[\gamma^\mu,\gamma^\sigma\gamma^\rho]
=
2g^{\mu\sigma}\gamma^\rho
-
2g^{\mu\rho}\gamma^\sigma
\]

\[
[\gamma^\mu,S^{\rho\sigma}]
=
\frac{i}{4}
\left(
2g^{\mu\rho}\gamma^\sigma
-
2g^{\mu\sigma}\gamma^\rho
-
2g^{\mu\sigma}\gamma^\rho
+
2g^{\mu\rho}\gamma^\sigma
\right)
\]

\[
[\gamma^\mu,S^{\rho\sigma}]
=
\frac{i}{4}
\left(
4g^{\mu\rho}\gamma^\sigma
-
4g^{\mu\sigma}\gamma^\rho
\right)
\]

\[
[\gamma^\mu,S^{\rho\sigma}]
=
i
\left(
g^{\mu\rho}\gamma^\sigma
-
g^{\mu\sigma}\gamma^\rho
\right)
\]

\[
(\mathcal J^{\rho\sigma})^\mu{}_\nu\gamma^\nu
\]

\[
(\mathcal J^{\rho\sigma})^\mu{}_\nu\gamma^\nu
=
[\gamma^\mu,S^{\rho\sigma}]
\]


\subsubsection{}
Using the inf. expressions for $\Lambda_{1/2}$ at first order we have:
\begin{equation}
\begin{split}
    &\Lambda^{-1}_{1/2} \gamma^\mu  \Lambda_{1/2} = \left(I+\frac{i}{2}\omega_{\rho\sigma}S^{\rho\sigma} \right)\gamma^\mu \left(I-\frac{i}{2}\omega_{\rho\sigma}S^{\rho\sigma} \right) \simeq \gamma^\mu + \frac{i}{2}\omega_{\rho\sigma}[S^{\rho\sigma},\gamma^\mu] \\ &= \gamma^\mu - \frac{i}{2}\omega_{\rho\sigma} (\mathcal J ^{\rho\sigma})^\mu_{\ \nu} \gamma^\nu = \left[\delta^\mu_{\ \nu} - \frac{i}{2}\omega_{\rho\sigma} (\mathcal J ^{\rho\sigma})^\mu_{\ \nu} \right]\gamma^\nu \simeq (e^{-\frac{i}{2}\omega_{\rho\sigma}J{\rho\sigma}})^\mu_{\ \nu}\gamma^\nu = \Lambda^\mu_{\ \nu}\gamma^\nu
\end{split}
\end{equation}


\subsubsection{}
We denote $\Lambda \equiv \Lambda_{1/2}$, because only $\psi$ and $\bar\psi$ are transformed as we saw in (4):
\begin{equation}
    \begin{split}
        \bar \psi \gamma^\mu \psi \to (\bar\psi\Lambda^{-1}) \gamma^\mu (\Lambda\psi)=\bar\psi\Lambda^\mu_{\ \nu}\gamma^\mu \psi
    \end{split}
\end{equation}
since the defined lorentz transform $\Lambda^\mu_{\ \nu}$ acts trivially on $\psi,\bar\psi$ we can pull it out: $$\bar \psi \gamma^\mu \psi \to \Lambda^\mu_{\ \nu}\bar\psi \gamma^\mu \psi$$

\subsection{}
$\bar\psi$ has the conjugate velocity $\partial_\mu \bar\psi$ which doesn't appear in the Lagrangian leaving us with the first equation:
\begin{equation}
\begin{split}
    &\frac{\partial\lagr}{\partial\psi} - \partial_\mu\frac{\partial\lagr}{\partial(\partial_\mu\psi)} = 0\\
    &-m\bar\psi - \partial_\mu\left(i\bar\psi\gamma^\nu\delta^\mu{}_{\nu} \right) = 0\\
    &m\bar\psi + i\gamma^\nu\partial_\nu\bar\psi = 0
\end{split}
\end{equation}

\begin{equation}
\begin{split}
    &\frac{\partial\lagr}{\partial\bar\psi} - \partial_\mu\frac{\partial\lagr}{\partial(\partial_\mu\bar\psi)}=0\\
    & i\gamma^\nu\partial_\nu\psi - m\psi = 0
\end{split}
\end{equation}

\begin{equation}
    \boxed{\begin{cases} & m\bar\psi + i\gamma^\nu\partial_\nu\bar\psi = 0 \\ & i\gamma^\nu\partial_\nu\psi - m\psi = 0 
    \end{cases}}
\end{equation}




\end{document}